\begin{table}[htbp]\centering
\caption{Balance table: demographics across experimental conditions}
\label{tab:balance_demo}
\begin{threeparttable}
\begin{tabular}{lcccc}
\toprule
 & Pos Info & No Info & Neg Info & F-test (p) \\\\
\midrule
\multicolumn{5}{l}{\textbf{Panel A: Men}} \\\\
\addlinespace
Age & 22.12 (3.05) & 22.50 (2.65) & 23.47 (3.84) & 0.329 \\\\
Master & 0.12 (0.34) & 0.21 (0.42) & 0.05 (0.23) & 0.286 \\\\
Economics & 0.28 (0.46) & 0.39 (0.50) & 0.37 (0.50) & 0.645 \\\\
Number of experiments & 3.06 (2.29) & 3.68 (3.01) & 3.47 (4.40) & 0.744 \\\\
\addlinespace
\multicolumn{5}{l}{\textbf{Panel B: Women}} \\\\
\addlinespace
Age & 24.00 (5.21) & 22.15 (2.57) & 22.39 (3.80) & 0.268 \\\\
Master & 0.24 (0.44) & 0.12 (0.33) & 0.11 (0.31) & 0.448 \\\\
Economics & 0.35 (0.49) & 0.38 (0.50) & 0.14 (0.36) & 0.110 \\\\
Number of experiments & 3.12 (2.18) & 3.58 (2.72) & 3.36 (3.21) & 0.871 \\\\
\addlinespace
\bottomrule
\end{tabular}
\begin{tablenotes}\footnotesize
\item Notes: Cells report mean (SD) within gender and experimental condition. F-test reports the p-value from testing equality of means across the three conditions within each gender (regression with condition indicators; joint test).
\end{tablenotes}
\end{threeparttable}
\end{table}
